% This module uses amsmath, amssymb, amsthm, and the main document's
% definition, lemma, corollary, and proof environments.
% The normalization lemma below is a new, manually proved interface result.
% The separately cited Tamarin theorem is an existing symbolic result.

\section{Reservation origins, persistent records, and recovery}
\label{sec:security-model}

A write to one session can reseal a journal while preserving another
session's saved version and reserved randomness. The resulting image
belongs to the new write, but the untouched reservation still belongs
to its original draw. Describing this execution requires three identities:
one for the draw, one for the saved version, and one for the image that
carries it. The same distinction lets us follow a pending write before
installation and a copy acquired before its wrapping key.

The adversary schedules public messages, process loss, honest recovery
and the disclosures specified below for the asynchronous machine
$\mathcal P$. A successful capture gives the adversary the retained
bytes and leaves the authoritative journal in place. Adversarial rollback
or forking would require an explicit transition that replaces or branches
that history. Application deliveries remain private to the honest
applications. We take the protocol, its primitives and its public
outputs as parameters of the record-normalization lemma. With finite
budgets for execution, record sizes and queries, the wrapper's bookkeeping
is polynomial in their total representation size.
Table~\ref{tab:record-notation} gives the notation.

\begin{table}[t]
\centering
\caption{Principal objects in the record model. The three identities separate
an original random draw, an immutable session-state version and its physical
storage representation.}
\label{tab:record-notation}
\small
\begin{tabularx}{\linewidth}{@{}>{\raggedright\arraybackslash}p{.25\linewidth}X@{}}
\toprule
Notation & Meaning\\
\midrule
$o,\rho_o$ & Original invocation and its reserved random tape\\
$v,X_v$ & Immutable session-state version and its payload\\
$h$ & Physical image or pending object carrying one or more versions\\
$\kappa$ & Identity of actual key material, shared by equal keys\\
$\mathsf{res}(o)$ & Event at which the original random tape is drawn\\
$\mathsf{save}(v)$ & First retention of $v$ on the durable acquisition surface\\
$\mathcal C_t,\mathcal L_t$ & Acquired and readable original versions at prefix $t$\\
\bottomrule
\end{tabularx}
\end{table}

\subsection{Three origins and their event order}

\begin{definition}[Reservation origin]
\label{def:reservation-origin}
A reservation origin is a fresh invocation index $o$, together with
\[
  \bigl(\mathsf{kind}_o,\mathsf{owner}_o,s_o,e_o,B_o,
        \mathsf{res}(o)\bigr).
\]
Here $\mathsf{kind}_o$ distinguishes generation, encapsulation, and any other
reserved operation; $s_o,e_o$ identify its session and epoch; $B_o$ is the full
canonical binding, including the admitted policy and operation inputs; and
$\mathsf{res}(o)$ is the event at which its random tape $\rho_o$ is drawn.
Execution or exact replay of $o$ consumes the same $\rho_o$ and the same
binding.  A new draw receives a new origin even if a failed attempt's external
operation identifier is reused.
\end{definition}

The ideal source in this model draws independent tapes of the prescribed
lengths.  Explicit entropy-disclosure queries reveal designated tapes.
Independence is a distributional assumption about the source at sampling.
Later API calls preserve the sampled tapes. The generation and encapsulation budgets
$G,C$ count original draws, including draws in attempts that subsequently
fail.  Exact retries consume no new slot.  Independently sampled operations
retain distinct origins even if some public outputs happen to coincide.

A record version $v$ has an immutable payload $X_v$ and semantic metadata
\[
  \bigl(\mathsf{subject}(v),\mathsf{role}(v),
         \mathsf{epoch}(v),\mathsf{phase}(v)\bigr).
\]
The subject is the session whose state the version records. A sibling
session writing the enclosing journal preserves that subject. A single
version may contain several reservations and state derived from earlier
ones. Preparing a version gives it a handle. Its first durable retention
is a different event, $\mathsf{Saved}(v)$, at time $\mathsf{save}(v)$.
For example, a pending intent may already hold an intact target before
the journal installs it. A caller can receive an unknown-outcome error
after the storage operation has reached this point.

A physical object $h$ has a separate creation or sealing index.  It may carry
one version, several sessions' versions, an exact earlier image, or a pending
target.  An unchanged version keeps its $v$ and $\mathsf{save}(v)$ when it is
copied, archived, applied from a pending intent, or included in a newly sealed
aggregate.  A new semantic session-state version receives a new $v$ even when
some of its fields are unchanged.  These rules distinguish the original
reservation from each session-state update that retains it.

The following transition rules specify the provenance discipline.

\begin{enumerate}
  \item \emph{Reserve and execute.}  Reserve creates $o$ at the actual draw.
  Execution uses its retained binding and tape.  Unknown-outcome
  reconciliation keeps that origin; a definite new attempt allocates another.
  \item \emph{Prepare and retain.}  Preparing $v$ fixes $X_v$ and its provenance.
  Retaining the first intact image or pending target containing $v$ emits
  $\mathsf{Saved}(v)$.  An attempt that retains no such object has no successful
  acquisition event.
  \item \emph{Copy, apply, and reseal.}  An exact copy preserves the complete
  object and its sealing origin.  Resealing creates a new sealing invocation
  and its actual new wrapper randomness, while preserving the referenced
  record versions.  Applying an existing pending target copies its retained
  bytes.
  \item \emph{Crash and resume.}  Process loss removes volatile references.
  Recovery reconstructs the retained handles and their original provenance.
  It does not turn a retained reservation into a new draw.
  \item \emph{Confirm and delete.}  Confirmation changes protocol state and may
  authorize particular host-side deletions.  Neither action changes a
  reservation's origin or removes an adversary's previously acquired copy.
\end{enumerate}

For an original reservation actually retained in $v$, causality requires
\begin{equation}
  \mathsf{res}(o)\ \leq\ \mathsf{save}(v).
  \label{eq:origin-causality}
\end{equation}
A later physical object's creation time dates the representation carrying
the reservation. Both original events in~\eqref{eq:origin-causality} remain
fixed, so the interval between drawing $\rho_o$ and first retaining it stays
explicit. The
symbolic reservation-and-save transition can represent that interval only
through a separately justified source correspondence.

\subsection{Acquisition, delayed readability, and key identity}

Wrapping-key handles resolve to identities of actual key material.  Two
session labels that designate the same key resolve to the same identity
$\kappa$.  Disclosure through either label therefore enables every acquired
wrapper protected by $\kappa$.  A rotation creates a different identity when
it creates different key material.  Signing-key identities are treated in
the same way across sessions and roles.

For each intact physical object $h$, its construction records a finite set
\[
  \mathsf{paths}(h)=\{(v,\gamma)\},
\]
where $\gamma$ is the set of wrapping-key identities on a recoverable
serialization path to $X_v$.  Multiple paths to one version are permitted.
An unprotected record has $\gamma=\varnothing$.  Nesting a wrapper adds its
key identity; aggregation unions its members' paths.  An exact ciphertext
copy keeps its paths.  Resealing decoded records reconstructs the paths under
the new wrapper and retains the same $v$ values.  Thus resealing decoded
plaintext and wrapping an old ciphertext are distinct operations.

Only intact, recoverably embedded record payloads create these paths.  A
public hash or MAC of a private record is represented by its computation and
output in the renderer introduced below. Its use of private bytes creates a
public observation, while an acquisition path requires an intact embedded
payload. The model therefore records a public head query as a commitment
observation.

\begin{definition}[Record acquisition and readability]
\label{def:record-acquisition}
A successful $\mathsf{Capture}(h)$ returns the actual retained bytes of $h$
and adds an immutable copy to the adversary's archive $\mathcal A_t$.
Let $\mathcal K_t$ be the wrapping-key identities disclosed by time $t$.
Define
\begin{align}
  \mathcal C_t
    &=\{v:\exists h\in\mathcal A_t,\ \exists\gamma,\
                    (v,\gamma)\in\mathsf{paths}(h)\}, \label{eq:acquired-records}\\
  \mathcal L_t
    &=\{v:\exists h\in\mathcal A_t,\ \exists\gamma,\
          (v,\gamma)\in\mathsf{paths}(h)
          \ \wedge\ \gamma\subseteq\mathcal K_t\}.
          \label{eq:readable-records}
\end{align}
Each capture records $\mathsf{AcquireOriginal}(v,t)$ for every version in its
support, including versions already acquired through another path.
An unsuccessful request for unavailable state returns its ordinary error and
adds no acquisition.
\end{definition}

Suppose the adversary acquires its first copy of $v$ in an encrypted
image $h$ at $t_1$. Its only path to $v$ still lacks one wrapping key.
The host deletes $h$, and the key is disclosed at $t_2>t_1$. The copy
in the adversary's archive survives the deletion. Thus $v$ enters
$\mathcal C_{t_1}$ when the bytes arrive and $\mathcal L_{t_2}$ when
the key arrives. There was only one capture. Both sets describe the
declared disclosures and key identities. Guessing a value changes neither.

An additional key that is explicitly recoverable from an opened record can be
included by a declared extraction rule.  In this extension, iterate
\eqref{eq:readable-records} and the declared key-extraction rules to their least
fixed point after each event.  Each rule identifies its input bytes, its
deterministic extraction algorithm, and the resulting actual key identity.
Only rules with available inputs fire.  This accounts for nested protection
without assuming that inaccessible keys have already been learned.

Signer-file capture is a separately typed disclosure.  For correspondence
with the retained symbolic theorem, it records exposure of the actual signing
key conservatively at acquisition, as does that theorem's signer-file rule.
The readability ledger can additionally record when the file's contents
become decryptable.  In particular, a wrapping-key disclosure cannot be
scoped to traffic state when the declared capture also includes the matching
signer file.

\paragraph{Constructive exposure.}
From public protocol inputs, explicit leaks, and readable record payloads, we
form a labelled replay-dependency graph.  Its edges are specified algorithms:
unseal a reservation with its matching key; regenerate a key from its original
tape; replay encapsulation on the bound offer; decapsulate a published
ciphertext; derive a root; and derive the direction-specific traffic and ACK
material.  Directly stored roots, skipped keys, and plaintext enter through
their own record fields.  For the implemented root schedule, an edge is
\[
 (R_{e-1},S_e,C_e)\longmapsto
 R_e=\mathsf{HKDF}\bigl(R_{e-1},S_e,
       D_{\rm root}\Vert H_{\rm core}(C_e),32\bigr).
\]
Closing these edges over the bounded execution yields a computable set of
message keys, each accompanied by the operations that reconstruct it. These
are positive exposure certificates. For messages outside the declared
disclosure conditions, the confidentiality game below measures the adversary's
remaining advantage. Confirmation counters record public progress while the
graph retains the dependencies of each replay computation.

\subsection{An original-record implementation of the copy interface}

We can now ask whether each physical copy needs its own private payload.
The following comparison normalizes copies to the original records they
carry and concerns record management independently of KEM secrecy.
The following two storage machines use the same protocol machine
$\mathcal P$, random-tape source, scheduler, and cryptographic interfaces.

The \emph{literal-copy machine} materializes the immutable payload $X_v$ in
each indicated serialization position.  Its objects include current images,
pending targets, and retained historical copies.  It records the origin tags
alongside these values as private bookkeeping.

The \emph{original-record machine} keeps one immutable table entry
$\mathcal O[v]=X_v$ and represents each serialization position by a reference
to $v$.  Its wrapper maintains the physical-object directory, pending states,
adversary archive, origin tags, and key identities.  It has the following
specified interfaces.

\begin{enumerate}
  \item $\mathsf{RegisterOriginal}(v,X_v)$ is a private write by the honest
  machine.  It fixes one value at a fresh $v$.  The wrapper receives the
  handle and the provenance metadata.
  \item $\mathsf{RenderOriginal}(\theta)$ evaluates the supplied serialization
  recipe $\theta$ by dereferencing $\mathcal O$.  A recipe specifies literal
  public fields, record references, references to already rendered objects,
  and each sealing, hash, MAC, or other admitted byte-processing call.
  A fresh sealing call uses a fresh invocation index and the same prescribed
  coin distribution as the literal-copy machine.  An exact-object read or
  replay returns its cached original bytes.
  \item $\mathsf{AcquireOriginal}(v,t)$ records acquisition of that version.
  This is an event, not an assertion of decryption.  A capture's externally
  returned bytes come from the rendered object cache.
  \item $\mathsf{RevealOriginal}(v)$ returns the original payload $X_v$.
  The wrapper uses this interface for a plaintext record-opening response
  only after a captured path is enabled by the disclosed keys, and caches
  the answer for subsequent openings.
\end{enumerate}

The renderer in this comparison has the original payload table, so it can
compute ciphertexts, commitment digests and report MACs directly from the
private bytes. A KEM embedding changes the information available to the
reduction: the challenger withholds the selected tape. The next reduction
must preserve the record interface from the information it actually receives.

The operational contract has three further requirements.
First, all publicly visible payload-dependent bytes and statuses are obtained
through the specified renderer or through $\mathcal P$'s common protocol
computation.  Memory addresses and host timing are not additional outputs of
this interface.
Second, actual byte-operation order, error responses, partial transaction
outcomes, and cryptographic-oracle calls are retained.  Caching an exact
object does not discard a native hash-validation call, observable retry, or
write attempt: such an action remains a renderer or protocol action.
Third, the reference metadata and the key-identity relation are available to
the wrapper without inspecting a selected hidden reservation's coins.
These are checkable interface conditions; none asserts the secrecy of a
record, a root, or a message.

\begin{lemma}[Alias normalization]
\label{lem:alias-normalization}
For the operations and renderer contract above, every adversary
$\mathcal A$ interacting with the literal-copy machine has an online wrapper
$\mathcal B$ interacting with the original-record machine such that their
public views have the same distribution.  Their reservation-origin ledgers,
record-save events, acquisition events, and read sets are identical after
projection onto the original handles.

The wrapper makes at most $|\mathcal L_T|$ distinct
$\mathsf{RevealOriginal}$ calls by terminal time $T$.  Repeated physical
copies, captures, and openings do not multiply this number.  The wrapper
preserves every original reservation and every fresh sealing invocation, so
it does not merge their respective cryptographic resource counts.

The statement also holds for an adversary with quantum private state and
quantum access to common public cryptographic oracles, provided the storage
interface is classical and both experiments retain the same oracle calls.
\end{lemma}

\begin{proof}
Maintain a map from every literal record occurrence to its immutable handle
$v$, and from every native object identifier to a rendered-object entry.
Couple the protocol random tapes, sealing coins, scheduler choices, and public
cryptographic oracles between the experiments.  The invariant is that
dereferencing every original-record recipe produces the corresponding
literal-copy bytes, with the same semantic provenance and physical
disposition.

For a new record, register its payload once and extend the occurrence map.
For a fresh rendering, submit the same recipe and byte-operation sequence to
$\mathsf{RenderOriginal}$.  Dereferencing gives the same inputs to every
sealing, digest, MAC, validation, and output operation.  Coupled coins and
common oracles therefore give the same bytes, statuses, and state changes.
This step retains a write that became durable before returning an error;
the caller's status does not overwrite the persistence state.

An exact copy or pending-target application copies the object reference and
its original cache entry.  Aggregation collects the corresponding original
references, and resealing evaluates the new recipe using those references.
Each operation preserves the session-record handles and their save times.
It follows by induction on the recipe constructors that
$\mathsf{paths}(h)$ is the same in both machines.  A public hash or MAC
constructor contributes its common output and has an empty record-acquisition
path set.

On successful capture, output the cached object bytes in both experiments
and append its immutable descriptor to the adversary archive.  For each
version in its support, emit the same acquisition event at the actual
capture time.  Preserve the ordinary error for an unavailable object.
When a matching wrapping key is disclosed, use the common key identity and
reevaluate the enabled paths in the complete adversary archive.  For an
opening response, obtain $X_v$ from $\mathsf{RevealOriginal}$ on its first
use and from the wrapper's cache thereafter.  Every such call has
$v\in\mathcal L_T$.  Shared-key aliases and later key acquisition are thus
handled without another copy of $X_v$.  The optional key-extraction extension
preserves the invariant by induction on the least-fixed-point rounds.

Host deletion removes only host references.  The adversary archive and
already returned payloads persist.  Crash and resume use the same retained
object identifiers; neither operation recreates a reservation or changes a
save event.  Ordinary protocol actions are performed by the same
$\mathcal P$ with the same results.  This completes the induction, including
adaptive choices based on any preceding response.

For quantum private state, identical classical responses and the same public
oracle transformations induce the same transformation of that state at
each interaction.  The argument introduces no measurement of a quantum
oracle query and no oracle-dependent advice.  Finally, caching each
plaintext opening by $v$ proves the query bound; reservation and fresh-seal
indices were preserved separately throughout.
\end{proof}

The adversary may retain several ciphertexts for one version and observe
a different public head for each. Alias normalization preserves all of
these observations, and a single opening of the common payload serves
every alias of that version. Each fresh seal and each public hash or MAC
computation still incurs its own cost.

% Manuscript fragment. Requires amsmath, amssymb, amsthm.
\subsection{Explicit guarded normalization of record readability}
\label{sec:guarded-provenance}

One session version can contain several reserved values, and several
records can contain the same value. Suppose we are given the objects
and keys acquired up to time $t$. Which original draws can we now read?
A source $o\in\Omega_{\rm src}$ names one such draw, and
$\mathcal E^{\rm src}_t$ records the identities of the sources whose
protected values we can derive. We continue to use $\mathcal C_t$ and
$\mathcal L_t$ for versions and $\mathcal O[v]$ for the payload of $v$.

Here we fix the graph of admitted records and the public key-derivation
edges. Key derivation in this fragment is independent of recovered
plaintext. The preceding interface specifies how to add an explicit
plaintext-dependent key-extraction rule. The provenance rules supply
the identities in the graph, and the renderer supplies the corresponding
bytes. A second copy carries the same source identity and can give the
adversary another route to its value.

\paragraph{Immutable record graph.}
Let $\mathcal R$ be a finite set of immutable encrypted record objects.
Each $r\in\mathcal R$ has a matching decryption-key identity $\kappa(r)$,
a set of embedded child records $\operatorname{child}(r)$, and a set
$\operatorname{orig}(r)$ of protected source values in its plaintext.
Containment is acyclic.  Record identities distinguish ciphertext objects;
source identities distinguish original sampling operations.  Several
records can contain the same source.  A rewrapped record has a new record
identity and may have a new key, while retaining its source identity.
Key identities are equivalence classes of equal actual keys. Equal key bytes
therefore have one identity across separate API objects.

Let $\mathcal D\subseteq\mathcal K\times\mathcal K$ be fixed public key
derivation edges.  An edge $(k,k')$ means that possession of $k$ permits
deterministic computation of $k'$.  This interface represents, for example,
derivation of an operation-wrapping key from the journal key.  In the
theorem below, these edges do not depend on plaintext recovered from records.
Write $D^*$ for their reflexive transitive closure.

Let $\Omega_{\rm src}=\bigcup_{r\in\mathcal R}\operatorname{orig}(r)$ be the
declared source domain.  At a trace prefix $t$, let $A_t\subseteq\mathcal R$ be all record objects
ever acquired by that prefix, $K_t\subseteq\mathcal K$ all directly
disclosed keys, and $\mathcal S^{\rm dir}_t\subseteq\Omega_{\rm src}$ all directly disclosed source values, represented
by their source identities.  These are cumulative sets.  The decryption-key
closure is
\[
  K_t^*=\{k':\exists k\in K_t,\ (k,k')\in D^*\}.
\]
The concrete derivation rules are
\[
 \frac{k\quad(k,k')\in\mathcal D}{k'},\qquad
 \frac{r\quad\kappa(r)}{c}\ (c\in\operatorname{child}(r)),\qquad
 \frac{r\quad\kappa(r)}{o}\ (o\in\operatorname{orig}(r)).
\]
Their least closure starts from $A_t,K_t,\mathcal S^{\rm dir}_t$.  A fact $o$ in the source
sort means that the protected value of source $o$ is derivable; it does
not mean that a public identifier has become known.  The source projection
of this closure is denoted $\mathcal E^{\rm src}_t$.

\paragraph{Guarded routes.}
For every path $p=(r_0,\ldots,r_\ell)$ in the containment graph ending
at a record with $o\in\operatorname{orig}(r_\ell)$, define
\[
 g(p)=\{\kappa(r_0),\ldots,\kappa(r_\ell)\}.
\]
The normalized table contains a route $(r_0,g(p),o,p)$ for every such
path.  Its exposure expression is
\[
 \widehat{\mathcal E}^{\rm src}_t=\mathcal S^{\rm dir}_t\ \cup\
 \{o:\exists(r,g,o,p),\ r\in A_t\ \land\ g\subseteq K_t^*\}.
 \tag{GP}
\]
One available path to $o$ suffices. Along that path every required
key must be available, and a key used twice need only be known once.

\begin{theorem}[Exact guarded normalization]
\label{thm:guarded-normalization}
For every finite record graph and derivation interface above, and every
acquisition prefix $(A_t,K_t,\mathcal S^{\rm dir}_t)$, the concrete and normalized exposure
sets agree: $\mathcal E^{\rm src}_t=\widehat{\mathcal E}^{\rm src}_t$.  The equality applies to all prefixes of
an adaptively chosen history.
\end{theorem}
\begin{proof}
Key derivations depend only on disclosed keys and the fixed relation
$\mathcal D$, so their least closure is exactly $K_t^*$.
For the forward inclusion, take a source $o\in \mathcal E^{\rm src}_t$.
If it was directly disclosed then $o\in \mathcal S^{\rm dir}_t$.
Otherwise its first derivation decrypts a record $r_\ell$ with the
matching key.  If that record was not acquired directly, its first
availability follows by decrypting a parent.  Follow these first
derivations backwards.  The finite acyclic record graph makes this
process terminate at an acquired record $r_0$.  Every intervening key
lies in $K_t^*$, yielding a route satisfying (GP).
For the reverse inclusion, take a satisfied route.  Starting with its
acquired first record, successively decrypt along the path.  Every
required key is in $K_t^*$, so each child and finally $o$ is in the
concrete least closure.  Direct disclosures are in that closure by
definition.  Both inclusions hold for each fixed prefix; no independence
or advance choice of the acquisition schedule was used.
\end{proof}

\paragraph{Incremental evaluation.}
Equation (GP) answers the readability question by listing the paths to
each source. A small record DAG can have exponentially many such paths.
As acquisitions and disclosures arrive, we can instead track the facts
that each event makes derivable. The evaluator uses indexed forward
chaining for Horn rules~\cite[Section~4]{dowlingGallier1984}.
Write
\[
\begin{aligned}
 N&=|\mathcal R|+|\mathcal K|+|\Omega_{\rm src}|,\\
 M&=|\mathcal D|+\sum_{r\in\mathcal R}
       \bigl(1+|\operatorname{child}(r)|+|\operatorname{orig}(r)|\bigr).
\end{aligned}
\]
The unit in each summand counts the record's matching-key incidence.
In an adjacency-list representation, repeated incidences are counted with
multiplicity in $M$.

\begin{proposition}[Incremental exposure evaluation]
\label{prop:guarded-worklist}
For the fixed record and key-derivation graphs above, the exact exposure
set at every prefix of a cumulative acquisition and disclosure history
can be maintained with $O(N+M)$ preprocessing, $O(N+M+U)$ total time and
$O(N+M)$ working space, where $U$ counts input acquisition, key-disclosure
and direct-source-disclosure events, including repetitions. These bounds
use indexed atom identities with constant-time membership and insertion;
they include reporting each newly exposed source once. Exporting complete
sets at selected prefixes adds their total output size.
\end{proposition}
\begin{proof}
Index each key's outgoing derivation edges and the records having that
matching key. Maintain known keys, available records, opened records and
exposed sources. Enqueue a key or record only when it first becomes known
or available. A key event propagates its outgoing derivations and checks
its indexed records. A record event checks its matching key. When both
the record and its key are present, open the record once, making its
children available and its sources exposed. Direct disclosures insert
their sources immediately. Drain the queue after each input event.

Every insertion is an instance of a stated derivation rule, so all
maintained facts belong to the least closure. At an empty queue, every
derivation edge with a known parent has been processed. For a record
whose two premises hold, the later arrival of those premises triggers
its opening unless it was already opened. Thus every enabled rule has
fired, and the maintained facts are exactly the least closure. Applying
this argument after each input event and using
Theorem~\ref{thm:guarded-normalization} gives the exposure set (GP) at
every prefix, including adaptively chosen prefixes.

Each key and record is queued at most once. Each key-derivation and
matching-key incidence is scanned at most once from its key, each
record's key is tested once from its record, and each opened record's
payload incidences are processed once. Repeated input events take
constant time. The total work is therefore $O(N+M+U)$, and the indexes,
sets and queue occupy $O(N+M)$ space. Each source enters the incremental
output at most once; complete-set exports are charged separately.
\end{proof}

The graphs are fixed throughout this bound. Local deletion changes future
acquisition availability and leaves the cumulative closure intact. A new
record or derivation edge changes the indexed input and requires its own
update accounting. An implementation using dictionaries has the
corresponding expected bounds under constant expected lookup and insertion
costs; normalization of external identities is charged separately.
The finite checker retains explicit route enumeration
for its small fixtures, where it compares the guarded normal form with
forward decryption. The indexed evaluator avoids that enumeration;
materializing all route witnesses can itself require exponential output.

\begin{corollary}[Copies, deletion and late disclosure]
\label{cor:guarded-history}
Exposure is monotone under additional acquisitions, direct disclosures
and key disclosures. Copying or rewrapping a source preserves existing
routes and may add routes to the same source identity.
Local deletion of a record previously acquired
does not remove its routes from adversarial knowledge.  In particular,
a later matching-key disclosure can expose a source in a ciphertext
acquired before that disclosure and deleted locally in the meantime.
\end{corollary}
\begin{proof}
The cumulative sets and their key closure only grow, and every predicate
in (GP) is positive. An exact ciphertext copy reuses its record identity
and route. Source-preserving rewrapping may change record identities or
key guards while retaining the terminal source. Local deletion
restricts future acquisition from that local store; it removes no member
of $A_t$ and therefore cannot invalidate a satisfied route.
\end{proof}

\paragraph{What a readable source tells us.}
For a source $o$, a path from an acquired record and the keys that
open every wrapper on that path give a witness for its readability.
Direct disclosure enters $o$ in $\mathcal E^{\rm src}_t$ immediately.
Absence from the set has a precise meaning within this fragment. These
record decryption and fixed public key-derivation rules have no way
to derive the value from the history so far. Other ways of learning a
secret, including guessing, cryptanalysis, timing leakage or access to
a quantum decryption oracle, require their own model. For a claim
about the selected application message we simulate the native
ciphertexts and commitments and bound the adversary's advantage in
the protocol experiment.

\paragraph{Native shape and finite checks.}
At the fixed research source, operation reservations are encrypted under
a wrapping key derived from the journal key, and tokens are retained
within encrypted durable images or pending targets.  The corresponding
route therefore includes the image key and the token key, with their
actual derivation edge.  A fresh aggregate image can preserve an older
token and its original source.  Public image-head digests are distinct
observations from the record objects needed by the decryption rules.
The executable supplement independently evaluates iterative decryption
and guarded-route normalization, exhausts all acquisition/key/direct-source
subsets for its declared finite graphs, and checks histories with late
key disclosure and local deletion.  Its counts and source hash are in
the accompanying machine-readable result; those checks validate the
implemented finite model and supplement the general proof above.

% Standalone manuscript fragment; no preamble.
% Requires proposition and proof environments (e.g. amsthm).
\paragraph{Why exposure must retain conjunctive guards.}
Could we label each item with the sources it exposes and take the
union of those labels as items arrive? A single encrypted record and
its wrapping key settle the question.

\begin{proposition}[Acquisition exposure is not additive]
\label{prop:acquisition-nonadditivity}
Let $r$ be a retained-record atom, $k$ a wrapping-key atom, and $s$ a
protected-seed atom associated with an origin identifier $o$.
For $A\subseteq\{r,k\}$, let $\operatorname{Cl}(A)$ be the least set
containing $A$ and closed under the rule
\[
  \{r,k\}\;\Longrightarrow\;s.
\]
There are no other rules in this fragment.  Define the exposed-origin set by
\[
  E(A)=\{\,o : s\in\operatorname{Cl}(A)\,\}.
\]
There is no fixed function
$f:\{r,k\}\longrightarrow\mathcal{P}(\{o\})$ satisfying
\[
  E(A)=\bigcup_{x\in A} f(x)
  \qquad\text{for every }A\subseteq\{r,k\}.
\]
\end{proposition}

\begin{proof}
The closure rule gives
\[
  E(\{r\})=E(\{k\})=\varnothing,
  \qquad
  E(\{r,k\})=\{o\}.
\]
Exactness for the two singleton acquisitions forces
$f(r)=f(k)=\varnothing$.  Their union is then empty for the joint
acquisition, contradicting $E(\{r,k\})=\{o\}$.
\end{proof}

\paragraph{Meaning of exposure.}
The origin identifier $o$ may be public. Membership in $E(A)$ says that
the rules let us derive its protected seed. To see how combining
observations can change an estimator's success probability, take two
independent uniform $\lambda$-bit strings $S,W$ and let
$C=S\mathbin{\oplus}W$. Assume there is no further observation that
depends on either string. Then
\[
  H(S\mid C)=H(S\mid W)=\lambda,
  \qquad H(S\mid C,W)=0.
\]
From either observation alone, any estimator of $S$ succeeds with
probability at most $2^{-\lambda}$; from both, $S=C\mathbin{\oplus}W$
recovers it exactly. This XOR construction illustrates acquisition in
information-theoretic terms. The cryptographic record interface handles
the protocol's actual record encryption.

\paragraph{Consequence for provenance routes.}
A fixed origin label identifies the seed carried by a copy. Readability
also depends on how the observations combine: the record and its key
may arrive separately and enable recovery only together. Provenance
routes retain this conjunction. Copies retain the same source identity,
and each route becomes usable when its record and every required key
are available.


\subsection{A recovery experiment with acquisition-based eligibility}

\begin{definition}[Fixed-target recovery experiment]
\label{def:recovery-game}
Fix a protocol profile $\mathcal P$ and finite execution budgets, together with
a designated session $s$, recovery cut $c$, target epoch $e$, and candidate
generation and encapsulation reservation handles $g_e,a_e$.  Handles may be
named by original invocation slots before those invocations occur; a missing
invocation will fail the reservation-event test below.  Public
protocol messages may be delayed, duplicated, reordered, or replaced.
The adversary can capture the declared retained objects, reveal wrapping
keys, reveal actual signing keys or signer files, and reveal designated
reservation tapes.  The cut exposes the selected wrapping keys; preceding
roots may already be public or disclosed by the stated prior captures.
These disclosures have their exact scope in the event ledger.

For one fresh logical application send, the adversary supplies equal-length
messages $(m_0,m_1)$.  The challenger samples $b\leftarrow\{0,1\}$ and runs
the ordinary sending transition on $m_b$.  Rejection or an unknown outcome
keeps its ordinary protocol behavior.  A retained accepted send is identified
by $(s,e,d,j)$, including direction $d$ and the epoch-scoped logical index $j$;
exact retries reuse its original ciphertext.  The challenge is valid only
when that logical send is accepted in the designated epoch.  Subsequent
protocol actions and disclosures continue until the adversary outputs $b'$.
\end{definition}

Write $\iota_A,\iota_B$ for the two actual signing-key identities configured
for the target session.  The predicate
$\mathsf{ReservedFor}(s,e,g_e,a_e)$ checks the occurrence of the named honest
generation and encapsulation reservation events in the target epoch's local
durable plans, with their target session and roles.  It checks event metadata,
without asserting that their outputs form the accepted target exchange.
Eligibility uses the following event conditions over the complete execution:
\begin{align}
 \mathsf{Eligible}_{s,e,c}(\tau)\ \Longleftrightarrow\quad&
 \mathsf{ValidTarget}(\tau)
 \ \wedge\ \mathsf{ReservedFor}(s,e,g_e,a_e)
 \nonumber\\[-1mm]
 &{}\wedge\ c<\mathsf{res}(g_e)
          \ \wedge\ c<\mathsf{res}(a_e)
 \nonumber\\
 &{}\wedge\ \iota_A,\iota_B
                  \notin\mathsf{IdentityExposed}(\tau)
 \nonumber\\
 &{}\wedge\ \mathsf{EntropyExposed}(s,e;\tau)=\varnothing
 \nonumber\\
 &{}\wedge\ \nexists v\in\mathcal C_T:
       \mathsf{subject}(v)=s\ \wedge\ c<\mathsf{save}(v).
 \label{eq:eligibility}
\end{align}
Here $\mathsf{EntropyExposed}(s,e;\tau)$ records the target epoch's designated
entropy disclosures.  The separate event $\mathsf{Attr}$ states that these
named reservation events form the accepted target's unique usable
generation/encapsulation pair, with an adapter-computable association to the
original slots.  In particular, it connects the generation public key, the
encapsulation input, the authenticated response, and the target's subsequent
key schedule.  This matching is a
source/authentication obligation: a failure of $\mathsf{Attr}$ can occur while
$\mathsf{Eligible}_{s,e,c}$ holds and is charged in the computational
accounting.

Consider a journal image sealed after $c$ by a write to another session.
What does acquiring that image mean for the target session $s$? Read the
last line of~\eqref{eq:eligibility} using the first save of each version
of $s$ in the image. If all of them were first saved by $c$, this
capture adds no post-cut version of $s$. A single version first
saved after $c$ violates it, even while the bytes are encrypted. The
test uses $\mathcal C_T$ across all roles and epochs of $s$.
$\mathcal L_T$ answers when the acquired versions become readable.
Other epochs' entropy disclosures keep their usual behavior, and a send
may precede confirmation by both peers.

Define the unconditional centered advantage by
\begin{equation}
 \mathsf{Adv}^{\rm rec}_{\mathcal P}(\mathcal A)
   =
 \left|
   \mathbb E\!\left[
    \mathbf 1_{\mathsf{Eligible}_{s,e,c}(\tau)}
                       (-1)^{b\oplus b'}\right]
 \right|.
 \label{eq:recovery-score}
\end{equation}
Invalid or ineligible executions contribute zero. Keeping these executions
in the probability space lets the reduction compare unconditional
expectations, including failures and adaptive disclosures. Its score thus
remains defined even when the adversary changes the probability of eligibility.
If every execution with $\mathsf{Eligible}_{s,e,c}=0$ is completed by an
independent fair guess, the conventional advantage over one half is
\[
 \mathsf{Adv}_{\rm msg}
   =\left|\Pr[\text{completed guess is correct}]-\tfrac12\right|
   =\tfrac12\,\mathsf{Adv}^{\rm rec}_{\mathcal P}(\mathcal A).
\]
The computational accounting uses this latter convention.

The adaptive experiment still has one application challenge. After
interacting, the selector can name its coordinates and reservation
handles, provided they identify the same accepted logical send. The
test bit stays with that send. For each fixed pair in $[G]\times[C]$,
the reduction scores executions in which the selector attributes the
send to that pair. It evaluates eligibility, disclosures and attribution
from each comparison world's own complete history, allowing the
selection to depend on the transcript.

The record restriction in~\eqref{eq:eligibility} is the one used by the
retained multi-epoch symbolic theorem: its \textsf{Snapshot} action outputs
a protected record, while \textsf{Cut} separately releases the wrapping keys.
The theorem establishes symbolic plaintext secrecy with both original
reservations after the cut, key-scoped identity exclusions, the target
entropy exclusion, and the stated post-cut saved-record acquisition
exclusion~\cite{qperiaptSymbolic2026}.  Its joint replay verifies the main
lemma and all 24 reusable dependencies.  Equation~\eqref{eq:recovery-score}
specifies the computational message experiment to which the cryptographic
modules are to be connected; the alias lemma establishes its record-interface
normalization.

\begin{corollary}[Safe stopping at a forbidden acquisition]
\label{cor:target-acquisition-stop}
Consider the fixed-target experiment of
Definition~\ref{def:recovery-game}.  Suppose an actual successful capture
contains a target-session version $v$ with $c<\mathsf{save}(v)$, identified
by the original-record metadata.  For any terminal bit $D$, replacing the
remainder of that execution by output zero leaves
\[
   \Pr[\mathsf{Eligible}_{s,e,c}(\tau)\wedge D=1]
\]
unchanged.  Replacing its terminal score by zero also leaves the expectation
inside~\eqref{eq:recovery-score} unchanged.

In particular, if a captured target-session version retains the selected
reservation $o\in\{g_e,a_e\}$, then
$c<\mathsf{res}(o)$ and~\eqref{eq:origin-causality} establish this stopping
condition.  The wrapping key may be obtained before or after the capture.
\end{corollary}

\begin{proof}
The successful capture inserts $v$ into $\mathcal C_t$ immediately.
Its save time and subject are immutable, and acquisition remains recorded
after host deletion, resealing, confirmation, or process loss.  Thus the
last conjunct of~\eqref{eq:eligibility} is false for every continuation.
Each such continuation already contributes zero to both quantities.
On executions that never reach this condition, the experiment is unchanged.
The final assertion follows from
$c<\mathsf{res}(o)\leq\mathsf{save}(v)$.
\end{proof}

The successful acquisition event and its record metadata determine the
stopping test before a challenge adapter serializes the response's private
bytes. The reduction stops because every continuation already has zero
eligible score; the native experiment retains the SDK's capture behavior.
A public digest, an unrelated record, or an unsuccessful capture keeps its
real behavior, so permitted digests and report MACs must still be supplied
through their declared interfaces.
For adaptive target selection, the same reasoning applies within each
fixed-candidate comparison. The stopping certificate applies to that
candidate; the original game may still select another eligible target.

The same timing issue arises for the reservation itself. A generation
tape reserved before $c$ can still be the tape used after a delayed
receipt. In that execution the reservation condition
in~\eqref{eq:eligibility} fails even if the KEM computation completes
later. The two endpoints' confirmed counters record their own progress
during the delay.

\subsection{Connection to the cryptographic argument}

The normalization theorem fixes the original objects whose bytes, acquisition,
and key identities the reduction must represent. The next section gives
constructive image/intent and report interfaces for a selected KEM operation
whose public outputs are available while its original random tape is opaque.
Their algorithms preserve original-object caches, public commitments, exact
retries, and actual durable outcomes. The fixed-target stopping rule above
extends to an adaptive provenance certificate in
Lemma~\ref{lem:adaptive-origin-stop}.

Usable attribution is computed from available semantic metadata. Original
generation and encapsulation slots belong to a fixed lifetime budget
$[G]\times[C]$: a fresh draw receives a new slot, while copies retain their
source. Each candidate induces its own protected records and report-key
closure. Proposition~\ref{prop:native-local-accounting} decomposes the score
in the native experiment before comparing those local continuations. It
retains both unattributable eligible executions and permitted key projections
as explicit native-game terms.

The computational composition evaluates accepted operations, eligibility, and
all public outputs anew in each comparison world. Exact component repeats
retain one response and one correctness event. The native constructor map in
Section~\ref{sec:native-correspondence} identifies the storage operations used
by these interfaces; Section~\ref{sec:discussion} specifies the joint root and
message continuation to which their bounds connect.
